\chapter*{Preface}
\addcontentsline{toc}{chapter}{Preface}

The RSVP framework begins from a deliberate reversal of the usual order of exposition in physical theory. Where standard treatments derive algebra from an already-given geometry, the present volume derives geometry as a shadow cast by algebra. The plenum is not first a manifold upon which fields are later defined; it is first an algebra of interacting operators, and geometry emerges only once that algebra is asked to describe its own smooth deformations. Curvature, on this reading, precedes space rather than living within it.

This ordering is not a notational preference. It reflects a substantive claim: that the primitive fact about reality is misalignment between potential, flow, and entropy, and that everything ordinarily called curvature, from the bending of light to the bending of a life, is a record of that misalignment. The present book develops the algebra adequate to this claim before any metric structure is introduced, so that later volumes in the series can treat geometry, category, and sheaf as successive refinements of a single algebraic fact rather than as independent postulates.

The scalar potential $\PhiField$, the vector flow $\vField$, and the entropy $\SEntropy$ together form the RSVP triple $\RSVPTriple$. Each field is continuous; none of the constructions in this book depend on discretization, and the reader should resist the temptation to think of the plenum as built from discrete units later smoothed by approximation. The smoothness is constitutive, not emergent from a coarser substrate.

The book may be read in two registers at once. Read as mathematics, it is a study of commutators, ideals, and cohomology on a specific algebra of fields. Read as philosophy, it is an account of how persistence, purpose, and obligation arise from the same non-commutativity that, mathematically, gives rise to curvature. Neither register is metaphorical with respect to the other; the claim of the RSVP program is that they are the same content viewed from two vantage points, and the chapters below move between them without marking the transition, since no transition, properly speaking, occurs.

